\documentclass[../main.tex]{subfiles}
\begin{document}
\section{Overview and scope}
\label{sec:overview}

This document is a single, unified treatment of three functional inequalities --- the Poincar\'e
($\CP$), log-Sobolev ($\CLS$), and quadratic transportation-cost ($\CTCI$) constants --- for
probability measures $\Pstar(\dd x)\propto e^{-U(x)}\dd x$ that arise in statistics and machine
learning, together with the deepest structural question the subject raises, the
Kannan--Lov\'asz--Simonovits (KLS) conjecture. It is organized as a \emph{difficulty gradation}, in
three parts.

\paragraph{Part I --- structured statistical and ML families.}
For concrete, parameterized families the answer is reachable now by routing the family to a
classical proof template. The reading order mirrors the increasing difficulty:
\begin{itemize}[leftmargin=1.4em]
  \item Section~\ref{sec:def} fixes the three constants and the normalization;
  \item Section~\ref{sec:results} recalls the imported toolkit (Bakry--\'Emery,
        Brascamp--Lieb, Holley--Stroock, tensorization, Otto--Villani, Hardy) and maps the
        known landscape across common distributions;
  \item Sections~\ref{sec:tier1}--\ref{sec:tier5} climb the curriculum --- Gaussian and
        log-concave calibration (Tier 1), bounded perturbations and products (Tier 2), the
        Gaussian-prior GLM-posterior spine (Tier 3, the flagship), heavy-tailed targets and
        the obstructions they trigger (Tier 4), and multimodal / hierarchical posteriors
        (Tier 5);
  \item Section~\ref{sec:agenda} collects the open research agenda.
\end{itemize}

\paragraph{Part II --- A-series results and open-target deep dives.}
The research agenda of Section~\ref{sec:agenda} is then developed in depth, one section per
flagship problem: data-informed constants for GLM posteriors (A1), the Bernstein--von Mises
Gaussian limit (A2), heavy-tailed posteriors beyond classical LSI (A3), variational-inference
transport constants (A4), and the quotient / reparameterization of symmetry- and
funnel-induced slow modes (A5). These sections now contain twenty-one independently reviewed results and
sharp calibrations as well as open problems. Each separates proved or imported results, useful
obstructions, analytic drafts pending review, and the residual theorem-level targets.

\paragraph{Part III --- the KLS frontier.}
The dimension-free Poincar\'e bound for an \emph{arbitrary} isotropic log-concave measure is the KLS
conjecture \cite{KannanLovaszSimonovits1995,LeeVempala2018}, open since 1995. The published
benchmark is $\CP\le C\,(\log D)\,\norm{\Cov}_\op$ \cite{Klartag2023Logarithmic}; a July
2026 version-1 preprint improves the claimed bound to
$\CP\le C\sqrt{\log D}\,\norm{\Cov}_\op$ \cite{Letwin2026QuadraticKLS}. It is the ``Tier-$\infty$''
boundary of Part I: every structured family there buys its explicit constant precisely by having
more usable structure than a general log-concave measure. Part III is organized in six groups. It
opens with an orientation to the conjecture, the August 2026 frontier, and the covariance-spike
obstruction that constrains every route (Section~\ref{sec:kls-orientation}), followed by a
self-contained survey of the five strategy families in the literature
(Sections~\ref{sec:family-needles}--\ref{sec:family-transport}). It then develops this
repository's own work: the shared localization machinery
(Sections~\ref{sec:notation}--\ref{sec:models}), the conditional fixed-cut Eldan program
(Sections~\ref{sec:introduction}--\ref{sec:open}), the fixed-eigenfunction spectral route
(Section~\ref{sec:spectral-route}), and the distinct deterministic moment-map/Haar route with its
commutator bottleneck (Section~\ref{sec:moment-map-cmh}). That last route carries a second,
logically prior layer (Sections~\ref{sec:cmh-normalization}--\ref{sec:cmh-exact-cases}) in which
the proposed endpoint is given precise operator data, proved to dominate the affine Poincar\'e
constant, and computed exactly on the line, on products, and on every log-concave Dirichlet law
--- together with the two reasons that endpoint may nonetheless be false. It closes by mapping each remaining gap
in the literature onto an open node of the KLS ledger (Section~\ref{sec:kls-synthesis}).

\paragraph{The guiding intuition.}
A single picture organizes Part I. Strong curvature of $U$ pulls mass back to equilibrium
and makes all three constants small ($\Hess U\succeq mI\Rightarrow\CP,\CLS,\CTCI\le 1/m$);
flat directions, heavy tails, and separated modes are the three ways this fails ---
producing, respectively, infinite or large constants, the loss of LSI/$T_2$, and
exponentially large constants. The Gaussian $N(0,\Sigma)$ is the reference point against
which every deformation is measured.
\end{document}
